\chapter{稠密地图表示}

\textbf{作者：} Victor Reijgwart、Jens Behley、Teresa Vidal-Calleja、Helen Oleynikova、Lionel Ott、Cyrill Stachniss 和 Ayoung Kim

本章将关注 SLAM 的另一个方面，即建图组件。我们假设机器人位姿已经完成估计，目标是构建其周围环境的稠密地图。典型方法首先使用 SLAM 后端（前几章已讨论）求解机器人的运动轨迹，然后在给定轨迹的情况下重建稠密地图。本章将详细介绍稠密地图表示，重点讨论地图元素、数据结构和构建方法。

早期建图方法主要采用第 1 章所述的稀疏、基于路标的方案，仅从环境中提取少量显著特征。随着计算能力提升，以及机械式 3D LiDAR 或 RGB-D 相机等高精度 3D 范围传感器的出现（这些传感器能够以高频率提供详细的 3D 范围测量），稠密地图表示逐渐受到研究关注。对于规划、导航、操作和精确定位等需要详细理解环境的下游任务，稠密地图至关重要。本章解释这些稠密方法如何利用范围传感器数据的完整信息来细化和更新综合地图。

本章首先在第 5.1 节介绍支持稠密建图的主要传感器类型，重点是能够产生详细范围测量的范围传感器。第 5.2 节介绍基本表示元素和数据结构，并在第 5.4 节针对具体应用展开。与稀疏的基于路标建图相比，稠密地图表示的选择取决于所使用的传感器类型和预期的下游应用。影响表示类型选择的关键因素将在第 5.5 节总结。

\section{范围感知基础}

在深入讨论稠密地图表示之前，我们先简要总结一种 SLAM 中常用的关键感知模态，即范围传感器，为后续讨论提供必要背景。这类传感器测量环境中物体的距离，包括 LiDAR 传感器、飞行时间（time-of-flight，TOF）相机、RGB-D 相机和双目相机。这里我们重点介绍最常用的 LiDAR 传感器和 RGB-D 相机，它们主要用于室内外环境中的 SLAM 与稠密建图。

\begin{figure}[htbp]
\centering
\includegraphics[width=\linewidth,height=0.78\textheight,keepaspectratio]{397bf3ea6ff858f62e430dd83330be86e339a5e5f9cca5799262d9a8119765dd.jpg}
\caption{图 5.1 LiDAR 传感器的单射线和多射线类型。图中给出了真实场景中的 RGB、深度和 LiDAR 范围图像示例。}
\label{fig:chapter-5-397bf3ea6ff858f62e430dd83330be86e339a5e5f9cca5799262d9a8119765dd}
\end{figure}

\section{传感器测量模型}

先简要回顾感知机制及其测量模型。对于 LiDAR 传感器，\textsuperscript{1} 激光束由传感器发射，经环境反射后被探测器接收 \textsuperscript{[935]}。设激光束发射时刻为 $t _ { \mathrm{emit} }$，探测时刻为 $t _ { \mathrm{detect} }$，利用光速 $c$ 可按下式得到范围 $r$：

\begin{equation}
r = \frac {c (t _ {\mathrm{detect}} - t _ {\mathrm{emit}})}{2}.\tag{5.1}
\end{equation}

单条射线测量可以通过按特定模式移动的射线阵列扩展为二维（2D）或三维（3D）点集合，例如进行 360 度旋转或沿特定形状扫描。传感器生成的点集合称为点云，它是构建地图的基本元素。LiDAR 测量也可以表示为范围图像 $\pmb { R } \in \mathbb { R } ^ { B \times M }$：其中存储传感器完成一次完整旋转（即 360°）时 $B$ 条射线各自的范围。因此，对于每条射线，在传感器的水平视场内有 $M$ 个测量值；图 5.1 给出了示意和示例。

机器人应用中另一类常用范围传感器是 RGB-D 相机，例如 Microsoft Kinect、Azure Kinect DK 和 Intel RealSense。RGB-D 相机除了提供高为 $H$、宽为 $W$ 的 RGB 图像 $I _ { \mathrm {RGB} } \in \mathbb { R } ^ { 3 \times H \times W }$ 外，还提供一个尺寸相同的深度图 $\pmb { I } _ { D } \in \mathbb { R } ^ { H \times W }$，其中每个像素位置包含深度或范围。早期 RGB-D 相机使用具有已知图案的结构化红外（IR）光源，将该图案投射到环境中来生成深度图 $I _ { D }$；随后根据图案的畸变确定各像素的距离。由于阳光通常会干扰这种感知机制，采用投射光的 RGB-D 相机主要用于室内环境。新一代 RGB-D 传感器引入 IR 纹理投影器或飞行时间（TOF）技术，受干扰影响较小，因此可以支持室外应用。

\section{转换为点云}

利用范围传感器（例如 LiDAR 或 RGB-D 相机）的内参，可以将范围图像 $R$ 或深度图 $I _ { D }$ 转换为点云 $\mathsf { P } = \{ p _ { 1 } , \hdots , p _ { N } \}$。其中点 $\pmb { p } _ { i } \in \mathbb { R } ^ { 3 }$ 在传感器的局部坐标系中表示。点是地图表示中最基本的单元，本章后续将进一步讨论。

对于 LiDAR，传感器为每条射线提供内参标定 $( \phi _ { i , j } , \theta _ { i , j } )$，其中 $1 \leq i < B$ 且 $1 \leq j < M$；它们分别是方位角 $\phi _ { i , j } \in [ 0 , 2 \pi ]$ 与极角（俯仰角）$\theta _ { i , j } \in [ - \pi , \pi ]$，如图 5.2 所示。利用每条射线的已知角度，可以将 $R _ { i , j }$ 中的范围测量 $r _ { i , j }$ 转换为三维点 $\pmb { p } = ( x , y , z )$：

\begin{equation}
x = r _ {i, j} \cos (\theta_ {i, j}) \cos (\phi_ {i, j})\tag{5.2}
\end{equation}

\begin{figure}[htbp]
\centering
\includegraphics[width=\linewidth,height=0.78\textheight,keepaspectratio]{4cb8d131afc13808e90aab558038aa1c2b1707e944026628de3cf58d141cbe72.jpg}
\caption{图 5.2 球面坐标。}
\label{fig:chapter-5-4cb8d131afc13808e90aab558038aa1c2b1707e944026628de3cf58d141cbe72}
\end{figure}

\begin{equation}
y = r _ {i, j} \cos (\theta_ {i, j}) \sin (\phi_ {i, j})\tag{5.3}
\end{equation}

\begin{equation}
z = r _ {i, j} \sin (\theta_ {i, j}).\tag{5.4}
\end{equation}

对于 RGB-D 相机，通常采用针孔相机模型，将像素位置 $(u,v)$ 处的范围 $r _ {u,v}=R _ {u,v}$ 转换为三维坐标。利用相机内参 $\pmb { K } \in \mathbb { R } ^ { 3 \times 4 }$，将图像坐标中的齐次坐标 $\pmb { x } = ( u , v , r _ { u , v } )$ 通过伪逆转换到相机坐标中的点 $\pmb {p}$：

\begin{equation}
\pmb {p} = \pmb {K} ^ {\dagger} \pmb {x}.\tag{5.5}
\end{equation}

根据点的组织方式，所得点云可分为有组织点云和无组织点云。从深度图转换得到的点云是有组织的，每个点的位置与对应深度图的像素位置保持结构关系，因此可通过简单索引计算相邻点。LiDAR 传感器生成的点云则更复杂。例如，静止的 3D 机械式 LiDAR 生成的点云具有组织结构；而对于非重复图案固态或 flash LiDAR，以及运动中的机械旋转式 LiDAR，这种组织结构不再成立。很多情况下，无组织点云会因传感器运动而发生畸变，范围图像中的测量邻域也不一定与空间邻域相关。因此，需要估计传感器在完成一次扫描期间的运动（例如利用惯性测量单元（IMU）测量或里程计信息），结合每条射线的时间信息，对 LiDAR 点云进行去畸变，以补偿传感器运动 \textsuperscript{[1133, 256, 1231]}。关于运动畸变及其补偿的更多细节，请参阅第 8 章。

\begin{figure}[htbp]
\centering
\includegraphics[width=\linewidth,height=0.78\textheight,keepaspectratio]{f9bdfadb8891b58ab43430913880377c4f3a68a5a0aec8024ec3acf694c2cc0c.jpg}
\caption{图 5.3 常见稠密表示的示例。}
\label{fig:chapter-5-f9bdfadb8891b58ab43430913880377c4f3a68a5a0aec8024ec3acf694c2cc0c}
\end{figure}

\section{稠密建图基础}

地图是利用传感器信息和数据处理方法生成的、用于建模环境的符号结构 \textsuperscript{[1091, 136]}。需要注意的是，地图表示具有多样性，同一空间信息可以有许多不同表示方式（如图 5.3 所示）。场景表示的选择和精度会显著影响具体任务的性能，因此表示方式应由使用场景决定。例如，机器人中的运动估计和定位倾向于采用稀疏表示，如 3D 点特征 \textsuperscript{[787, 896]}，以利用这些特征实现一致的机器人位姿估计。另一方面，场景重建的关键目标是获得准确、稠密且高分辨率的地图，例如用于检查任务 \textsuperscript{[506, 822, 922]}。类似地，路径规划需要稠密信息，例如障碍物占据状态或与障碍物的最近距离 \textsuperscript{[774, 304, 1110]}。总体而言，本章围绕以下三个问题组织内容。

\textbf{问题 1：稠密建图需要估计什么量？} 表示环境时最常用的量是占据状态和距离。占据状态是建图中的关键属性，用于区分自由空间与被占据空间。距离估计则通过测量到邻近表面或物体的范围，为自由和占据空间提供更符合机器人视角的解释。

稠密建图基础。
\begin{figure}[htbp]
\centering
\includegraphics[width=\linewidth,height=0.78\textheight,keepaspectratio]{4462b836d740234d9efb94fe9df2eec98699164de4853c7ce956d811fcc039f2.jpg}
\caption{图 5.4 表示可以是显式的，也可以是隐式的。为清晰起见，图示已作简化；在实际中，抽象并非严格可分，通常可以组合使用。}
\label{fig:chapter-5-4462b836d740234d9efb94fe9df2eec98699164de4853c7ce956d811fcc039f2}
\end{figure}

\textbf{问题 2：应如何表示世界？} 这是关于采用何种空间抽象进行表示的问题。概括而言，表示可以是显式的，也可以是隐式的。显式空间抽象还可根据所使用的几何类型分类，而隐式表示可根据所选函数分类。图 5.4 展示了这些抽象类型。

\textbf{问题 3：应使用什么数据结构和存储方式？} 选定的表示需要存储在内存中供后续使用。数据结构和存储方法应根据具体应用及预期用途进行选择。

文献中存在大量利用范围传感器生成场景稠密表示的方法，它们在估计量、空间抽象、存储结构、连续性和应用领域等方面各不相同。

在讨论不同表示之前，我们先研究由范围测量估计的主要量，重点理解建图阶段最常估计的量。下面简要概述每个量的基本定义，并说明其重要性及发挥作用的具体场景。

\section{占据地图}

占据栅格由 Elfes 和 Moravec \textsuperscript{[304, 774]} 在三十多年前提出，此后得到广泛使用。占据栅格\textsuperscript{2} 简单且计算高效，因此在室内（甚至室外）环境建图中很受欢迎。最简单的情形是估计某个栅格单元被占据的概率：将单元含有障碍物建模为概率，已占据单元的占据概率为 1，被判定为空的单元为 0。本质上，占据建图是一个二元分类问题，用于预测每个单元属于二元类别的概率。

\begin{figure}[htbp]
\centering
\includegraphics[width=\linewidth,height=0.78\textheight,keepaspectratio]{a0c3917803db17ea2e6794fb0653ac66b0c83186ab12ac2eb9393a3bd34af8dc.jpg}
\caption{图 5.5 二维简单射线投射示例。给定某个（计算或估计的）方位角处的范围测量，范围测量的回波表示存在障碍物。}
\label{fig:chapter-5-a0c3917803db17ea2e6794fb0653ac66b0c83186ab12ac2eb9393a3bd34af8dc}
\end{figure}

给定传感器测量集合 $z _ { 1:t }$ 和传感器位姿集合 $\pmb { x } _ { 1:t }$，地图 $m$ 中每个单元被占据的概率可表示为 $p ( \pmb { m } | \boldsymbol { z } _ { 1:t } , \pmb { x } _ { 1:t } )$。图 5.5 给出了示例地图。假设每个单元 $m _ { k }$ 相互独立且测量条件独立，占据状态可用著名的对数几率形式高效更新 \textsuperscript{[774]}：

\begin{equation}
l (m _ {k} | \pmb {z} _ {1: t}) = l (m _ {k} | \pmb {z} _ {1: t - 1}) + l (m _ {k} | z _ {t}),\tag{5.6}
\end{equation}

其中 $l (\cdot | \cdot ) = \log ( o (\cdot | \cdot ) )$，$o (\cdot | \cdot )$ 为几率形式：

\begin{equation}
o (m _ {k} | \boldsymbol {z} _ {1: t}) = \frac {p (m _ {k} | \boldsymbol {z} _ {1 : t})}{1 - p (m _ {k} | \boldsymbol {z} _ {1 : t})}.\tag{5.7}
\end{equation}

占据建图的主要优势体现在式（5.6）中：只需使用先前占据值和逆传感器模型 $l (m _ { k } | z _ { t } )$，通过简单加法即可更新概率。尽管如此，占据栅格要实现高效运行，依赖于对环境作出的很强假设。特别是，假设一个单元的占据似然与其他单元无关，会忽略空间相关性，而空间相关性对于推断附近未观测区域的占据状态可能十分重要。此外，传统占据栅格要求预先定义环境离散化方式，使地图各处空间分辨率保持不变。

\section{隐式表面与距离场}

表示机器人周围几何形状的另一种方法不是使用占据概率，而是描述自由空间与占据空间之间的边界。在理想情况下，可以用关于三个变量 $x,y,z$ 的解析函数描述该边界的形状和位置；当点 $\pmb { p } = (x,y,z)$ 位于表面上时，函数值为 0。换言之，该函数的零交叉对应表面本身。这种函数称为隐式表面。按照约定，当 $\pmb {p}$ 位于物体内部时函数符号为负，在外部时为正。图 5.6 给出了一个简单示例，说明如何确定隐式函数值。

\begin{figure}[htbp]
\centering
\includegraphics[width=\linewidth,height=0.78\textheight,keepaspectratio]{2dc1fb5d9d56f86eadd81aeb50c13cde2fb71fc3f6bf6d80e48c7596ee06da56.jpg}
\caption{图 5.6 用显式表面（黑色轮廓）和隐式表面（彩色场）表示的实心二维圆盘。隐式表面函数在物体内部（蓝色）为负，在外部（红色）为正。注意该函数的零交叉对应表面本身。}
\label{fig:chapter-5-2dc1fb5d9d56f86eadd81aeb50c13cde2fb71fc3f6bf6d80e48c7596ee06da56}
\end{figure}

连续隐式表面表示有许多优点。由于不存在固定分辨率，它可以在任意细节层次表示任意形状的物体。此外，只需计算函数符号，就能判断空间中的给定点位于障碍物内部还是外部。

可以使用不同类型的函数对隐式表面建模，其中欧氏有符号距离函数（Euclidean Signed Distance Function，ESDF）是一种常见选择。在任意查询点处，ESDF 的幅值表示到最近表面的距离，符号表示该点是否位于障碍物内部。ESDF 表示常用于加速几何算法，例如碰撞检查。此外，还可以从 ESDF 推导高质量的邻近梯度，用于基于优化的运动规划和形状配准。

ESDF 通过为地图中的每个点寻找（欧氏）最近表面点来计算。对于已知表面，可使用 Fast Marching 等技术高效完成。然而，在表面估计中，更常使用投影有符号距离，因为它可以从测量中高效计算，也更适合滤波。给定一条穿过查询点的测量射线，投影距离定义为沿射线从射线端点到查询点的距离，因此不必显式搜索最近点。尽管投影距离会高估欧氏距离，但其零交叉（估计表面）仍然正确。Curless 和 Levoy \textsuperscript{[234]} 提出的、后来由机器人领域 KinectFusion \textsuperscript{[801]} 推广的标准隐式表面估计方法，将所有测量的投影有符号距离进行加权平均。为减小高估的影响，投影有符号距离函数通常被截断到固定范围内，称为截断带，此时称为截断有符号距离函数（Truncated Signed Distance Function，TSDF）。注意，ESDF 可以从 TSDF 和占据地图高效计算，详见第 5.4.4 节。

\section{占据地图还是距离场？}

占据表示和隐式表示都是体积方法，它们的共同点是：通过在观测体积内处处估计感兴趣的量来建模几何形状。然而，两种表示的优先目标不同，哪一种更好取决于具体应用。下面简要总结两个关键差异。

\textbf{建模直接性：} 测量射线可以直接告诉我们空间的哪些部分是自由、被占据或未观测的，因此基于占据概率的地图更新所需的启发式方法和假设更少。相比之下，隐式表面通常对到表面的距离建模。对于已知表面，该距离可以精确计算；但对于部分测量，只能依赖前述 TSDF 等距离代理。

\textbf{平滑性：} 隐式表面地图天然比占据地图平滑，后者建模的是二元属性。隐式表面的平滑性带来许多好处，最重要的是使其可微，从而产生对许多应用有价值的邻近梯度。平滑性还能降低离散化导致的近似误差，并通过插值获得良好的亚像素分辨率估计。然而，由于不连续性无法平滑表示，隐式表面往往会漏掉细薄障碍物。

\section{地图表示}

\section{目标空间结构的显式性}

如图 5.4 所示，可以根据表示的显式性及其目标空间进行分类。在 3D 建图中，表示体积很直接；但在机器人建图中，表面同样重要，因为它支持各种下游任务。我们可以考虑四类主要表示：显式表面、隐式表面、显式体积和隐式体积表示。

对于表面，可以采用显式或隐式方式表示。显式表示将表面视为二维流形，旨在用场景中物体的边界刻画空间。能够表示边界的最简单抽象就是范围传感器生成的点云。另一种通用表面表示是多边形网格（第 5.4.3 节），由顶点、边和面组成。网格通过形成相连的闭合多边形（更常见的是三角形）来编码体积的定向表面。Surfel（第 5.4.2 节）也是常用抽象，广泛用于建图。表面表示是可视化应用的基础，也用于渲染仿真环境、增强现实、计算机辅助设计和 3D 打印。

3D 体积建模采用类似策略。基于点的朴素表示常用于 LiDAR SLAM；此外，基于占据或距离的体素（第 5.4.4 节）也是显式表示的常见选择。存储体积地图时，需要仔细考虑数据存储，以降低计算成本。体积建图也使用隐式表示，通常通过函数实现。高斯过程（Gaussian process，GP，第 5.4.5 节）和 Hilbert map（第 5.4.6 节）是隐式表示的典型例子。

\section{空间抽象类型}

\section{点}

给定范围测量，一种直接的稠密地图表示是使用点云。构建时，我们利用估计的全局位姿 $\pmb { T } _ { t } ^ { 1 }$，将时刻 $t$ 在局部坐标系 $\mathcal { F } ^ { t }$ 中记录的点云 $\mathsf { P } _ { t }$ 累积到全局地图点云 $\mathsf { P } _ { M }$ 中。通常假设 $\mathsf { P } _ { M }$ 的全局坐标系取为第一个点云的坐标系 ${ \mathcal { F } } ^ { 1 }$。

简单累积点云 $\mathsf { P } _ { 1 },\ldots,\mathsf { P } _ { t }$ 无法扩展到大型环境，因此常见策略是丢弃同一空间位置的冗余测量。为此，大多数方法 \textsuperscript{[1269, 256, 1133]} 使用高效的最近邻搜索（例如体素栅格或基于层次树的表示，见第 5.3.3 节）对点云进行下采样和存储，例如每个体素只存储有限数量的点 \textsuperscript{[256, 1133]}，或仅保存满足特定准则的点 \textsuperscript{[1269]}。也可以只保存关键帧，即显式存储少量点云，但这需要确定何时生成关键帧或子地图。

点云是最基本的表示形式，可以转换为其他表示，例如通过 Poisson 表面重建 \textsuperscript{[1131]} 转换为网格，或通过 Marching Cubes \textsuperscript{[1132]} 转换为有符号距离函数（Signed Distance Function，SDF）。但这需要额外数据，例如点的观测视点；将多次测量合并为点云地图时，该信息可能丢失。因此，通常需要对表面方向作出假设，以区分物体内部和外部区域。

\begin{figure}[htbp]
\centering
\includegraphics[width=\linewidth,height=0.78\textheight,keepaspectratio]{729c87cd13c54d9aa29a946444ace672caafb8ea761d9117908e37c7b16390c7.jpg}
\label{fig:chapter-5-729c87cd13c54d9aa29a946444ace672caafb8ea761d9117908e37c7b16390c7}
\end{figure}

\begin{figure}[htbp]
\centering
\includegraphics[width=\linewidth,height=0.78\textheight,keepaspectratio]{e06e97e692e44fa9c671f1e277e96c47c65f233bb1c254cc891cdcae60f35651.jpg}
\label{fig:chapter-5-e06e97e692e44fa9c671f1e277e96c47c65f233bb1c254cc891cdcae60f35651}
\end{figure}

(a)
\begin{figure}[htbp]
\centering
\includegraphics[width=\linewidth,height=0.78\textheight,keepaspectratio]{e447a927343799a60338d24906b1fd2ccff69f11421e6172382b7553e7a9cdb4.jpg}
\label{fig:chapter-5-e447a927343799a60338d24906b1fd2ccff69f11421e6172382b7553e7a9cdb4}
\end{figure}

(b)
图 5.7 对 KITTI 数据集 \textsuperscript{[373]} 第 07 序列中一组 LiDAR 扫描，通过累积点云和 surfel 生成的地图进行定性比较。(a) 点云地图，点的亮度表示 LiDAR 测量的反射率。(b) 基于圆盘的对应 surfel 地图。累积所有点云的完整地图占用 2.95 GB，而 SuMa \textsuperscript{[66]} 的对应 surfel 地图仅占用 160 MB。

\section{Surfel}

点云地图直接表示测量，但存储的点不包含表面信息，也无法表示点是从哪个方向测得的。通过 surfel \textsuperscript{[870]}，可以向点添加方向信息来编码这些内容。Surfel 通常表示为圆盘或椭圆盘 \textsuperscript{[549, 66, 1184, 104, 103]}，更一般地也可以表示为用高斯模型描述的椭球 \textsuperscript{[1041, 1042, 282, 1318]}。Splatting \textsuperscript{[1318, 103]} 不仅可以整合纹理信息，还能将重叠 surfel 混合为特定视点下连贯的渲染结果；更详细的讨论见第 14 章。

一个常用的圆形 surfel（表示圆盘）由位置 $\pmb { p } \in \mathbb { R } ^ { 3 }$、法向方向 $\mathbf { \boldsymbol { n } } \in \mathbb { R } ^ { 3 }$ 和半径 $r \in \mathbb { R }$ 定义。利用现代图形处理器（GPU）的能力可以高效渲染这些几何基元，从而加速点到 surfel 的关联并显著减少内存占用。

图 5.7 对比了基于点云和基于 surfel 的地图表示。稠密点云能够以很高的细节准确表示环境，但代价是占用大量内存。相反，将多次测量聚合为单个 surfel 会损失细节，但能在保留大尺度表面主要结构的同时显著降低内存使用量。

与使用 surfel（即小型圆形表面片元）显式表示表面几何密切相关的是正态分布变换（normal distributions transform，NDT）\textsuperscript{[82, 1035]}。NDT 将空间划分为体素，并用正态分布 $\mathcal { N } ( \mu , \Sigma )$ 近似体素内的点，其中均值 $\pmb { \mu }$ 和协方差 $\Sigma$ 由所含点估计。协方差的特征值 $\lambda _ { 1 } < \lambda _ { 2 } < \lambda _ { 3 }$ 及其对应特征向量 ${ \pmb v } _ { 1 }, { \pmb v } _ { 2 }, { \pmb v } _ { 3 }$ 可用于估计体素内的表面属性。对于平面表面（$\lambda _ { 1 } \ll \lambda _ { 2 }$），最小特征值 $\lambda _ { 1 }$ 对应的特征向量 ${ \pmb v } _ { 1 }$ 即表面法向。因此，对平面表面而言，NDT 表示一个 surfel，同时也能更准确地表示无法用单个 surfel 近似的点分布。从这个意义上说，NDT 是一种混合表示：空间被划分为体素栅格，因此具有显式性；而体素由正态分布表示，能够连续表示体素内部空间，因此又具有隐式性。

\section{网格}

点云和 surfel 地图虽然描述了局部表面属性，但由于没有建模表面连通性，仍然相对稀疏。为了获得更完整的表示，可以使用网格，将表面描述为一组相互连接形成多边形集合的点。这样就能表示水密表面、查询并插值新的表面点，并沿相连表面高效遍历。

在网格术语中，每个多边形称为一个面，每个角点称为一个顶点。最常见的是三角形网格（每个面由三个顶点围成）和四边形网格（每个面由四个顶点围成）。注意，任意多边形（面）都可以拆分为等价的三角形集合；因此，三角形网格不仅最简单，也是最通用的网格类型。

网格是一种非常灵活且节省内存的表示，因为面和顶点的数量可以直接适应表面复杂度及所需细节层次。例如，任意大小的平面只需两个三角面和四个顶点即可表示。此外，网格适合并行处理和渲染。网格常用于与计算机图形学交叉的应用，如渲染、表面分析、操作和形变，也广泛用于数字模型、仿真或基于表面的算法，例如地面机器人的路径规划。

\section{体素}

点云和网格适合表示沿表面定义的环境属性。然而，占据状态和有符号距离等估计量是在整个体积内定义的。一种直接存储和处理体积属性的方法，是在规则栅格上对其进行离散化。离散化的占据地图和有符号距离地图分别称为占据栅格和有符号距离场。

类比二维像素，3D 栅格中的单元称为体素。由于栅格结构规则，每个体素都可以分配唯一索引并存入数据结构。需要注意，体素只是一个容器，更形式地说是一个空间划分；其内容的含义因方法而异。在经典占据栅格中，体素值表示该体素内任意点被占据的可能性，因此地图任意点的占据值等于包含该点的体素值。但体素值不必表示一个常量小立方体。例如，有符号距离场中的体素值估计体素中心处的有符号距离；要获得任意点处的距离，需要查询该点邻近体素并通过插值得到其值。一些应用甚至使用（稀疏）体素栅格存储和高效查询非体积属性，例如点或表面颜色。

\section{连续函数}

函数是连续建图的关键抽象。在这种情况下，建图问题归结为拟合参数化或非参数化函数，即求解回归问题。上面提到的许多空间抽象都要求预先离散化环境，通常使地图各处空间分辨率固定；而连续函数（无论参数化还是非参数化）提供更大灵活性，可以重新计算分辨率，并通过插值填补数据缺口。

某些参数化函数，例如二维无穷直线 \textsuperscript{[1106]} 和三维平面 \textsuperscript{[533, 374, 1106]}，需要对环境作出严格假设，从而限制场景表示。但它们在内存和计算复杂度方面高效。基于控制点的函数（如 B-spline \textsuperscript{[932]}）或非参数化函数（如基于 GP 的函数）能够在更少假设下连续建模环境。从占据状态 \textsuperscript{[816]}、隐式表面 \textsuperscript{[1186, 636]}、距离场 \textsuperscript{[1195]} 到表面本身 \textsuperscript{[1122]}，尽管计算复杂度较高，基于 GP 的表示因其概率性质而成为表示环境的常用选择；概率性质可以对观测区域和未观测区域进行不确定性量化与推断 \textsuperscript{[378]}。

连续函数用于建图的一个关键优势是：如果选取至少一次可微的函数，就能提供梯度。梯度信息对于定位（计算表面法向）\textsuperscript{[1196]}、闭环检测（计算地形特征）\textsuperscript{[377]}、数据融合 \textsuperscript{[636]} 和规划 \textsuperscript{[1196]} 都很重要。

\begin{figure}[htbp]
\centering
\includegraphics[width=\linewidth,height=0.78\textheight,keepaspectratio]{bbd51a8abb86d14f4b9dbf9318ccd6d5f8876ce59b2e1abd8d16e42de8529bd1.jpg}
\caption{图 5.8 Marching Cubes 算法将立方体投影到隐式表面上，查询立方体顶点处函数值的符号，并查表确定这些值对应的 15 种配置之一。（图片来源：Ryoshoru，“Marching cubes”，CC BY 4.0 许可，来源：https://commons.wikimedia.org/wiki/File:MarchingCubesEdit.svg）}
\label{fig:chapter-5-bbd51a8abb86d14f4b9dbf9318ccd6d5f8876ce59b2e1abd8d16e42de8529bd1}
\end{figure}

\section{转换}

上述抽象并不总是单独使用；它们经常相互转换，或以多种形式同时使用。

例如，可以将点和网格等显式几何转换为隐式表面。一种灵活的距离计算方法是最近点查询：可以针对任意显式几何执行，并且只在需要的时间和位置按需计算。另一种方法是在规则栅格的所有点上使用波前传播计算有符号距离，这可以通过 Fast Marching 方法高效实现 \textsuperscript{[991]}。

反过来，将隐式表面转换为网格表示也很常见。将距离场转换为网格的原始技术称为 Marching Cubes \textsuperscript{[696]}。该算法将隐式表面划分到固定大小的立方体栅格中，并独立处理每个立方体。每个立方体根据八个角点处的隐式表面值生成一组三角形元素（图 5.8），再通过线性插值细化顶点位置。包括 BIM（Building Information Modeling）和 CAD（Computer-aided Design）模型在内的网格，又可以采样生成点或 surfel。

有时还需要将离散表示转换为连续表示。这通常通过对连续表示的参数求解优化问题来实现，使相对于离散数据的拟合误差最小。

\begin{figure}[htbp]
\centering
\includegraphics[width=\linewidth,height=0.78\textheight,keepaspectratio]{8ee90dccf0e08060ba9abfa39ab992f214f72150d351a5c34c0a1ab203b79f72.jpg}
\caption{图 5.9 (a) 通过函数 $f(\cdot)$ 将逻辑栅格坐标映射到朴素数组存储，函数将坐标唯一映射为线性索引。(b) 将无序数据直接存入数组结构。}
\label{fig:chapter-5-8ee90dccf0e08060ba9abfa39ab992f214f72150d351a5c34c0a1ab203b79f72}
\end{figure}

\section{数据结构与存储}

前几节介绍的抽象都需要存储在内存中。本节通过考察数据结构的选择及其优缺点，讨论各种抽象如何存储。

\section{朴素数据存储}

对于许多表示，简单的动态可调整大小数组是合理的起点。对于具有预定义空间划分的数据，需要确定要存储的数据类型，以及将空间坐标转换为空间索引坐标的函数。这常用于构建表示占据值或有符号距离值的地图。对于不规则数据，只需确定存储的数据类型，例如点云或 surfel。图 5.9 展示了使用映射函数存储有序数据以及存储无序点云的朴素方法。

这种朴素方法的优点是简单且支持快速随机访问。代价是表示可能需要大量内存。此外，读写和修改操作虽然很快，但改变表示的空间维度可能代价高昂，因为需要复制整个数据结构的内容。

\section{哈希表}

解决上述朴素存储局限的一种自然扩展是使用哈希表。该方法将地图划分为多个分片，对每个分片的坐标应用哈希函数得到单个值，并将分片数据存入以该哈希值为索引的表中。分片通常表示具有明确坐标的地图子区域，例如规则栅格中的立方体；这些立方体可以对应单个体素或固定大小的体素组（称为体素块），也可以存储各自子区域中的点、surfel 或网格片段。

哈希表保留固定数组的 $O(1)$ 查找时间，同时允许地图动态增长而无需重新分配。使用哈希表存储稠密地图时需要考虑三点：

1. \textbf{分片粒度：} 较小分片只在需要的位置分配数据，可以提高稀疏性。但分片数量不应过多，否则会降低哈希表的插入性能和内存效率。当哈希表存储仅沿表面存在的属性时，这一权衡尤为重要。
2. \textbf{哈希函数：} 理想的哈希函数应在表中均匀分布键，即使数据在空间上相邻（建图中经常如此）。
3. \textbf{冲突处理：} 处理哈希冲突的方法（例如链式法，即每个条目包含链表，或开放寻址）会显著影响哈希表性能。

在大多数情况下，哈希表在快速访问与高效插入、删除之间取得了良好平衡。但针对特定应用，通常需要初始调参才能获得良好性能。

\section{基于树的数据结构}

另一种只为环境中相关部分分配内存、同时高效存储空间数据的方式，是使用分层树结构表示。与哈希表类似，树通常支持高效访问和插入；其独特优势在于分层结构，可高效存储多分辨率数据，并加速最近邻搜索等空间操作。最常见的树变体包括 kD-tree \textsuperscript{[72]}、包围体层次结构（bounding volume hierarchy，BVH）\textsuperscript{[223]} 和八叉树 \textsuperscript{[752]}。

其中，八叉树支持高效邻域搜索和增量整合新测量。八叉树中的每个节点由一个称为八分体的子空间表示。八分体由中心 $\boldsymbol { c } \in \mathbb { R } ^ { 3 }$ 和边长 $e \in \mathbb { R }$ 定义，对应一个轴对齐包围盒。每个八分体最多有 8 个子八分体，其边长为 $\frac{1}{2}e$，如图 5.10 所示。

\begin{figure}[htbp]
\centering
\includegraphics[width=\linewidth,height=0.78\textheight,keepaspectratio]{62d2224194886c41e6b38d4a72042e4c1051e1b13121ba09342f5ad28a959c2f.jpg}
\caption{图 5.10 不同树层级的八叉树及其八分体示例。八叉树的每一层都将空间细分为更细粒度的八分体。注意，八叉树的更深层级只表示被占据的空间。}
\label{fig:chapter-5-62d2224194886c41e6b38d4a72042e4c1051e1b13121ba09342f5ad28a959c2f}
\end{figure}

通常只在叶八分体（即没有子节点的八分体）中存储点，并通过遍历树结构确定内部八分体中的点子集。

构建八叉树时，在包含点云 $\mathsf { P }$ 的轴对齐包围盒内迭代地将空间划分为八分体。每次划分都将 $\mathsf { P }$ 分成子集 $\mathsf { P } _ { 1 },\ldots,\mathsf { P } _ { 8 }$，对应范围为 $\frac{1}{2}e$ 的 8 个子八分体。非空子集 $\mathsf { P } _ { i }$ 形成中心为 $\boldsymbol {c}$、范围为 $\frac{1}{2}e$ 的子八分体；当达到指定八分体大小或最小点数时停止。构建完成后，通过遍历树结构并按需添加内部节点，可以高效执行更新和插入。当新数据落在根八分体之外时，可以创建新的根节点，并将新数据和旧根节点分配给其子节点，从而扩展树。

与体素栅格相比，八叉树只表示包含数据的子空间，因此能够高效存储被占据空间。但这种内存效率需要通过树遍历访问特定叶八分体，可能增加定位一个点所需的时间；同时必须显式表示树结构，带来额外内存开销。近期一些方法针对这些内存开销进行了改进 \textsuperscript{[305, 820, 70]}。

\section{混合数据结构}

为平衡内存需求与数据访问运行时间，一些数据结构以特定方式结合不同数据结构的优点，形成混合表示。在这种权衡中，我们接受较低的内存使用效率，以获得更高效的内存访问。

例如，哈希体素栅格 \textsuperscript{[808]} 将环境划分为固定大小的稠密块（如 $8\times8\times8$ 个体素），再把这些块存储在哈希表中，从而结合稠密体素栅格和哈希表的优点。借助哈希表的灵活性，只需为包含有意义信息（如靠近表面的位置）分配块；同时，在每个块内使用普通 3D 数组存储体素，使操作简单、高效，并适合 GPU 加速。

另一种选择是将哈希表与树结合。类似于哈希体素栅格，VDB\textsuperscript{3} 数据结构 \textsuperscript{[792, 791]} 将空间划分为哈希块，但在每个块内存储分层树。该结构提供层次树的全部优点，包括多分辨率表示和高效最近邻查询。然而，由于每个块大小固定，无论环境规模如何，树的最大高度都保持不变。因此，查找和插入可以在常数时间内完成，并显著快于纯树结构。

\section{构建地图：方法与实践}

\small
\setlength{\tabcolsep}{2pt}
\begin{center}
\begin{tabular}{p{0.317\linewidth}p{0.317\linewidth}p{0.317\linewidth}}
\toprule
节 & 空间抽象类型 & 表示的地图实体 \\
\midrule
5.4.1 & 点 & 表面 \\
5.4.2 & Surfel & 表面 \\
5.4.3 & 网格 & 表面（相连） \\
5.4.4 & 体素 & 占据状态或隐式表面 \\
5.4.5–5.4.6 & 连续函数 & 占据状态或隐式表面 \\
\bottomrule
\end{tabular}
\end{center}

表 5.1 所介绍建图方法的总结。

到目前为止，我们讨论了需要估计的目标量以及建图可用的各种空间抽象。本节将详细考察用于构建这些地图元素的方法。方法按其主要空间抽象分类，如表 5.1 所示。注意，一些方法会使用额外的空间抽象来提高性能，例如将点分组到体素中，以实现更高效的存储和更快的查询。

\section{点}

如第 5.3.2.1 节所述，简单累积测得点云来存储点无法扩展到大规模环境，还会造成测量冗余。因此，大多数方法 \textsuperscript{[1269, 256, 1133]} 将基于点的表示与体素栅格或八叉树结合，用于表示稠密地图。此外，数据结构的选择由高效最近邻搜索的需求驱动；最近邻搜索对于通过迭代最近点（iterative closest point，ICP）进行扫描配准至关重要，因为点对应关系需要迭代建立。

为了处理大规模环境，某些方法（例如著名的 LiDAR SLAM 方法 LOAM \textsuperscript{[1269]}）会滤波原始点云，提取角点和平面点，从而显著缩小点云规模。它们使用体素栅格只存储地图表示中的点子集，以剪枝冗余测量。基于 LOAM 中基于点的体素化，后续方法 \textsuperscript{[998, 1150, 839, 667, 894, 999]} 分别改进点提取 \textsuperscript{[998, 1150, 839]}、优化流程 \textsuperscript{[839, 667]}，或整合 IMU 信息 \textsuperscript{[894, 999]}。

另一类方法采用不同方式处理点云数据量，以避免依赖能力很强的特征提取方法。通过体素栅格 \textsuperscript{[256]} 对点云进行规则采样，可以显著减少每次 LiDAR 扫描的点数并去除潜在冗余信息。关键在于，体素栅格中的点并不聚合和平均，而是保留原始测量。沿着这一思路，Dellenbach 等人 \textsuperscript{[256]} 和 Vizzo 等人 \textsuperscript{[1133]} 使用该策略对输入点云下采样，在体素栅格地图中仅存储每个体素内受限数量的点。

总的来说，如第 5.3.2.1 节所述，（哈希）体素栅格具有双重作用：通过每个体素存储有限数量的点测量来抽象空间，并通过直接索引相邻体素加速最近邻搜索。

\section{Surfel}

Surfel 可以采用与点云类似的策略。值得注意的是，Stückler 和 Behnke \textsuperscript{[1041]} 以及 Droeschel 和 Behnke \textsuperscript{[282]} 的后续工作使用八叉树，在八叉树层级的多个层次表示 surfel，以进行数据关联。所谓多分辨率 surfel 地图通过累积的均值和协方差统计量间接表示 surfel，类似 NDT。

相比之下，Whelan 等人 \textsuperscript{[1184]} 将 surfel 存储为简单列表，并利用高效渲染技术生成投影，用于室内 RGB-D SLAM 的数据关联。在这种情况下，surfel 是显式几何基元，因此需要直接处理以相应更新 surfel 属性（即大小和方向）\textsuperscript{[549]}。Whelan 等人 \textsuperscript{[1184]} 的一项关键贡献是利用地图形变直接形变 surfel，而不是依赖 PGO；检测到闭环时，可以利用 surfel 表示的测量来形变地图。Park 等人 \textsuperscript{[846]} 也采用了类似的 surfel 地图形变策略。

同样，Behley 和 Stachniss \textsuperscript{[66]} 面向室外环境，需要用多个空间范围为 $100\,\mathrm{m}\times100\,\mathrm{m}$ 的子地图表示 surfel，并可将其从 GPU 内存换出。与 ElasticFusion \textsuperscript{[1184]} 不同，该方法依赖 PGO，但利用 surfel 可以自由定位这一点，将 surfel 绑定到位姿，从而可以用位姿图优化后的位姿直接形变地图；其他方法 \textsuperscript{[1158]} 也采用了这一策略。

\section{网格}

如前文所述，网格以富有表现力且灵活的方式表示相连表面。网格生成方法可分为两类。第一类直接将测得的点转换为网格；第二类将问题分为两步：重建隐式表面，然后提取等值面得到最终网格（见第 5.3.2.6 节）。

第一类方法通常直接计算输入点集的 Delaunay 三角剖分，并识别位于表面上的 Delaunay 三角形子集；\textsuperscript{[157]} 给出了详细综述。直接由点构建时，网格会隐式适应采样密度。这既可能是优点（提供自适应分辨率），也意味着方法对采样不规则和孔洞更加敏感。实践中，当整个表面可以由非常精确的深度传感器进行稠密采样时（例如使用测量设备），通常选择直接网格化方法。

第二类方法使用隐式表面作为中间步骤，在提取最终表面网格前简化数据融合和滤波。根据数据生成隐式表面的一种直观方法，是在规则栅格的每个点估计到表面的距离。如第 5.2.2 节所述，还必须根据每个点位于物体内部还是外部来设置隐式表面的符号。这一信息通常依据估计的表面法向确定，例如在小邻域上应用主成分分析（Principal Component Analysis，PCA）。然而，正如 \textsuperscript{[476]} 所指出的，这类方法可能产生不连续的隐式表面。为解决该问题，Carr 等人 \textsuperscript{[155]} 用一组径向基函数（Radial Basis Functions，RBF）对隐式表面建模，并通过求解全局优化问题拟合输入点。所得隐式表面在构造上是平滑的，并根据全局上下文可靠填补孔洞；但求解底层的大规模稠密优化问题计算代价很高。Shen 等人 \textsuperscript{[1002]} 使用移动最小二乘（moving least squares，MLS）对输入点进行局部近似，克服了这一限制。进一步地，Poisson 表面重建 \textsuperscript{[545]} 通过求解偏微分方程（partial differential equation，PDE）将隐式函数拟合到测量点的法向，得到稀疏且计算可处理的优化问题，对噪声尤其鲁棒。

在机器人应用中，通常希望从实时传感器流构建网格。通过增量更新，可以使表面重建达到实时效率。基于 TSDF 的表面重建因其天然的增量特性和总体简单性而广受实践采用。该方法属于第二类方法，通过平均投影距离估计隐式表面。由于更新 TSDF（或一般隐式表面）的代价远高于网格提取，实时方法主要集中于优化前者。

\section{体素}

基于体素的方法是 3D 重建和机器人学中最常用的体积表示之一。本节不按时间顺序罗列大量文献，而是聚焦实践中常见的概念，并依据三个基本决策标准组织内容：选择的估计量、数据结构和可扩展性。

\section{按估计量分类的方法}

基于体素的建图框架首先要选择估计量，最常见的选项是占据状态（见第 5.2.1 节）或距离度量（见第 5.2.2 节）；第 5.2.3 节的讨论可用于在两者之间作出选择。

自声呐的连续概率占据测量模型 \textsuperscript{[774]} 提出后，人们开发了简化的分段常数模型来降低计算成本 \textsuperscript{[478]}。这一转变受到 LiDAR 技术出现以及从 2D 地图向 3D 地图发展的研究兴趣推动。最近，Loop 等人 \textsuperscript{[692]} 提出连续概率模型：该模型不膨胀物体，而是使物体表面处的占据概率收敛到 0.5。占据估计由于召回率高而常用于碰撞规避，但与基于距离的方法相比，其不连续性和缺乏信息的梯度限制了性能（见第 5.2.3 节）。

对于距离度量，不仅要估计距离场的正值部分，还要在表面后方外推负距离，因为表面由有符号距离场的零交叉表示。为限制融合不完美正、负距离估计对精度的影响（见第 5.2.2 节），更新通常被限制在表面边界附近的小截断带内。然而，基于距离的方法仍容易擦除几何形状。例如，从相反两侧观测细薄物体时，对观测到的正距离和虚构的负距离求平均，会使零交叉翻转或消失。一些工作分析了截断带和权重衰减对最终重建质量的影响 \textsuperscript{[134]}；从根本上说，该问题可以减轻但无法消除。总体而言，在平滑表面上，TSDF 估计的表面优于占据方法，但对细薄物体的召回率较低。

TSDF 提供的距离信息本身非常有价值。然而，TSDF 并不保守，而是严格高估欧氏距离。为解决这一安全问题，voxblox \textsuperscript{[822]} 推广了增量构建 ESDF 的方法：它将传感器数据融合到 TSDF 中，然后使用 brushfire 算法 \textsuperscript{[629]} 更新 ESDF。随后，FIESTA \textsuperscript{[426]} 提出一种混合方法，直接从占据地图增量更新 ESDF。

\section{按数据结构分类的方法}

体积建图最简单的数据结构是静态 3D 数组。KinectFusion \textsuperscript{[801]} 表明，这种结构对于小型、固定大小的场景能够取得良好结果。但许多应用要求地图在运行时动态扩展，并且只在需要的位置分配体素以节省内存。

为此，Nießner 等人 \textsuperscript{[808]} 提出体素块哈希方案，将体素分组为块（例如 $8\times8\times8$ 个体素），并将块存入哈希表。该结构很快被 TSDF 采用，提供常数时间（$O(1)$）查找和动态插入；当然也可用于存储占据概率，如 FIESTA \textsuperscript{[426]} 所示。与单独哈希每个体素相比，将体素分组为块可以在哈希表大小和体素分配粒度之间提供可调权衡。

体素也可以使用树结构存储。Octomap \textsuperscript{[478]} 首先推广了使用八叉树存储占据概率，多年来一直是体积建图的事实标准。树结构的主要优点是天然支持多分辨率；主要限制是编码树结构会产生显著内存开销，且单元查找时间与树高成正比。近期方法通过混合数据结构缓解这一限制。例如，Supereight \textsuperscript{[1125]} 在前几层使用标准（动态）八叉树，在最后几层使用静态八叉树。静态八叉树可视为使用固定大小数组存储的八叉树，这消除了用指针编码父子关系的内存开销，但由于静态八叉树按块分配，粒度有所降低。VDB \textsuperscript{[792]} 数据结构最初为视觉特效（VFX）行业提出，后来被多个体积建图框架采用 \textsuperscript{[721, 1132]}。如第 5.3.3.4 节所述，它将块哈希与树结合，兼具两者优点：良好的内存效率、类似哈希表的常数时间查找和插入，以及类似树的多分辨率。

实际应用中，下游任务往往还要求在占据概率或 TSDF 旁存储颜色、语义 \textsuperscript{[398, 942]} 或 ESDF \textsuperscript{[822, 426]} 等附加信息。尽管几乎所有数据结构都可以扩展以支持附加通道，但实现工作量会随底层数据结构的复杂度增加。这进一步说明应优先使用简单数据结构（如体素块哈希）或灵活的第三方库。

\section{按测量融合算法分类的方法}

根据深度测量更新地图的算法称为测量融合器。它通过应用测量模型，更新每个观测体素的估计量。融合测量主要有两类方法：射线追踪和基于投影的方法。

对于每个测得的点，射线追踪融合器从传感器向该点发射射线，并更新射线穿过的所有体素。这种方法非常通用，只需知道传感器原点的位置。然而，一个体素可能被多条射线击中，尤其是靠近传感器的体素，这会造成重复计算；在并行实现中处理由此产生的竞争条件还会带来实现和性能开销。

相反，基于投影的方法直接遍历观测体素，并将每个体素投影到传感器坐标系，查询用于计算更新的射线。遍历地图而不是射线天然避免了竞争条件，因此基于投影的方法广泛用于多线程和 GPU 加速的体积建图框架。直接遍历地图产生的可预测访问模式还能减少内存瓶颈。但该方法的主要缺点是必须明确知道传感器的完整位姿和投影模型；对于无组织点云（包括 LiDAR 运动去畸变后的点云），该方法也更难使用。

\section{按可扩展性分类的方法}

内存和计算成本是体积建图的两大瓶颈。对于固定分辨率方法，内存和计算复杂度随地图总体积线性增长，并随所选分辨率的倒数立方增长。降低这些复杂度是重要研究方向，因为只有这样才能构建超越小型受限体积的高细节地图。

体积建图早期工作主要关注降低内存使用。例如，Octomap 使用八叉树内部节点存储子节点占据值的最大值或平均值。通过递归剪枝估计量接近父节点的叶节点，地图中的均匀区域会自动用更少、更低分辨率的节点表示。这种适应环境几何的策略在实践中非常有效，因为环境通常主要由自由空间组成。此外，在八叉树内部节点存储最小、最大或平均值对下游任务也有价值，可支持低分辨率地图查询，并使用层次算法执行快速碰撞检查或探索规划。但 Octomap 的核心限制是所有测量都以最高分辨率融合，因此其计算复杂度仍呈立方增长。

多分辨率还可用于降低测量更新的计算成本。由于测量射线以固定角度发射，击中远处几何体的射线更少，因此随着距离增加降低更新分辨率是合理的。可以通过多分辨率射线追踪 \textsuperscript{[286]} 或多分辨率投影融合 \textsuperscript{[1125]} 实现。Supereight2 \textsuperscript{[347]} 进一步根据测量更新的熵调节更新分辨率，从而降低计算复杂度。这些方法显著提高了更新性能，但不同分辨率层仍需显式同步。消除同步需求的一种方法是不在每个八叉树节点存储绝对值，而只编码各分辨率层之间的差异；形式上可通过小波分解实现。小波编码地图可以在任意时刻以任意分辨率高效查询。利用这一性质，wavemap \textsuperscript{[923]} 以从粗到细的方式更新地图，进一步降低计算复杂度。除根据测量熵调整更新分辨率外，它还跳过无信息更新，例如地图某区域的占据状态已经收敛为空闲且所有测量一致时。

\section{高斯过程}

如第 5.3.2.5 节所述，将建图问题表述为回归问题有利于获得连续表示。此外，若希望尽量减少对环境作出的假设，采用非参数方法求解非线性回归问题是理想选择。高斯过程（Gaussian process，GP）\textsuperscript{[909]} 是一种随机、非参数、非线性回归方法。给定其他位置处噪声且稀疏的测量，它可以估计任意查询点处未知函数的值。我们已在第 2.2 节介绍如何使用 GP 表示连续时间轨迹。显然，GP 也是连续量建图的有吸引力方案，并已广泛用于机器人学文献中，以利用深度传感器建模连续空间现象 \textsuperscript{[1122, 816, 377, 574]}。

GP 模型中的信息由均值 $\pmb { m }(\pmb { x })$ 和核函数 $\kappa(\pmb { x },\pmb { x } ^ {\prime})$ 给出，对所估计的连续量建模为：

\begin{equation}
f (\boldsymbol {x}) \sim \mathcal {GP} \left(\boldsymbol {m} (\boldsymbol {x}), \boldsymbol {\mathcal {K}} (\boldsymbol {x}, \boldsymbol {x} ^ {\prime})\right).\tag{5.8}
\end{equation}

设 $\mathsf { X }=\{\pmb { x } _ { j } \in \mathbb { R } ^ { D }\}$ 是具有测量值 $\pmb {y}$ 的位置集合，在位置 $\pmb {x}_j$ 处测得的估计量为 $y _ {j}=f(\pmb {x}_j)+\epsilon_j$。对于 $J$ 个训练对 $(\pmb {x}_j,y_j)$，假设噪声 $\epsilon_j$ 独立同分布，且满足高斯分布 $\epsilon_j\sim\mathcal {N}(\mathbf {0},\sigma_j^2)$。给定测试位置集合 $\mathsf { X } ^ {*}=\{\pmb {x}_n^*\in\mathbb {R}^D\mid n=0,\ldots,N\}$，函数值与观测目标值的联合分布可写为：

\begin{equation}
\left[ \begin{array}{c} \boldsymbol {y} \\ \boldsymbol {f} _ {*} \end{array} \right] = \mathcal {N} \left(\mathbf {1}, \left[ \begin{array}{c c} \boldsymbol {K} (\mathrm{X}, \mathrm{X}) + \sigma_ {j} ^ {2} \mathbf {I} & \boldsymbol {K} (\mathrm{X}, \mathrm{X} ^ {*}) \\ \boldsymbol {K} (\mathrm{X} ^ {*}, \mathrm{X}) & \boldsymbol {K} (\mathrm{X} ^ {*}, \mathrm{X} ^ {*}) \end{array} \right]\right),\tag{5.9}
\end{equation}

其中，$\pmb { K }=\left[\pmb {\mathcal { K }}(\pmb {x}_i,\pmb {x}_j)\right]_{ij}$。因此，条件分布 $(f_*\mid\mathsf {X},\pmb {y},\mathsf {X}^*)\sim\mathcal {N}(\overline{\pmb {f}}_*,\operatorname{cov}(\pmb {f}_*))$，其均值为：

\begin{equation}
\overline {{\boldsymbol {f} _ {*}}} = \boldsymbol {K} (\mathsf {X} ^ {*}, \mathsf {X}) [ \boldsymbol {K} (\mathsf {X}, \mathsf {X}) + \sigma_ {j} ^ {2} \mathbf {I} ] ^ {- 1} \boldsymbol {y},\tag{5.10}
\end{equation}

协方差为：

\begin{equation}
\operatorname{cov} \left(\boldsymbol {f} _ {*}\right) = \boldsymbol {K} \left(\mathrm{X} ^ {*}, \mathrm{X} ^ {*}\right) - \boldsymbol {K} \left(\mathrm{X} ^ {*}, \mathrm{X}\right) \left[ \boldsymbol {K} (\mathrm{X}, \mathrm{X}) + \sigma_ {j} ^ {2} \mathbf {I} \right] ^ {- 1} \boldsymbol {K} \left(\mathrm{X}, \mathrm{X} ^ {*}\right).\tag{5.11}
\end{equation}

式（5.10）和式（5.11）是所估计量的预测方程。GP 对表示具有空间相关性的数据尤其有效，因此能够克服环境建图中占据栅格方法所特有的单元独立性假设。高斯过程占据地图（Gaussian Process Occupancy Map，GPOM）\textsuperscript{[816]} 将传感器观测及其对应标签（自由或被占据）收集为训练数据；地图单元构成测试位置，并如式（5.9）所示与训练数据相关。利用式（5.10）和式（5.11）完成回归后，通过二元分类函数将回归输出“压缩”为单元的占据概率。

在原始形式中，GPOM 是一种批处理建图技术，计算复杂度为 $\mathcal {O}(J^3+J^2N)$。为应对这一计算复杂度，尤其是实现增量 GP 建图，后续提出了多种方法，例如 \textsuperscript{[574, 575, 1157, 378]}。

基于 GP 函数建图的一个关键优点是，可以对估计量进行线性运算 \textsuperscript{[969]}，结果仍为 GP。由于求导及梯度都是线性运算，估计量的微分输出具有概率性。环境不确定性的连续表示可用于突出未探索区域，并优化机器人的搜索规划 \textsuperscript{[378, 677]}。GP 地图的连续性也能提高规划器的灵活性，使其可以直接基于已收集的传感器数据推断，而不受栅格或体素单元分辨率限制。

\section{高斯过程隐式表面}

隐式表面也可以由 GP 表示。高斯过程隐式表面（Gaussian process implicit surface，GPIS）技术 \textsuperscript{[1186, 736, 677, 505]} 利用 GP 方法，根据噪声测量估计隐式表面的概率连续表示。此外，GPIS 也可连续地估计表面和距离场 \textsuperscript{[576, 1034, 636, 1195]}。

在 GPIS 表述中，设要估计的距离场为 $d$。给定表面上的点、这些点到最近表面的距离 $d_i$，以及通过线性算子计算得到的相应梯度 $\nabla \pmb {d}$ \textsuperscript{[969]}，由于表面上的距离为零，$\pmb {d}$ 与 $\nabla\pmb {d}$ 可以建模为零均值的联合 GP：

\begin{equation}
\left[ \begin{array}{c} \boldsymbol {d} \\ \nabla \boldsymbol {d} \end{array} \right] \sim \mathcal {GP} (\boldsymbol {0}, \boldsymbol {K} (\mathsf {X}, \mathsf {X} ^ {\prime})).\tag{5.12}
\end{equation}

GPIS 方法能够估计连续隐式表面及其由梯度得到的法向，并同时给出不确定性。一些工作考虑使用参数函数先验以更准确地捕捉给定形状 \textsuperscript{[736, 505]}；另一些方法不仅估计隐式表面，还估计完整距离场。就标准 GPIS 表述而言，距离在测量附近（即表面上）拟合良好，但在远离表面处会回归到均值，此处均值为零。为了在远离表面的位置以连续且概率的形式估计完整距离场，一些工作将非线性操作应用于类似 GPIS 的表述 \textsuperscript{[1195, 1196, 634]}。

所有这些方法都必须处理基于 GP 表述的计算复杂度；作为交换，仅凭点云就能获得环境的连续、生成式和概率表示。

\section{Hilbert map}

Hilbert map（HM）\textsuperscript{[904]} 在许多方面与 GPOM \textsuperscript{[816]} 相似。二者都是连续概率模型，与基于体素的方法不同，不对空间离散化；与基于点的方法相比，能够插值缺失数据。如前所述，GPOM 的主要挑战是计算开销高。因此，Hilbert map 的设计目标是：(i) 以在线方式连续处理数据；(ii) 建模观测之间的依赖；(iii) 纳入测量不确定性。

为实现这些目标，通常在投影特征空间中使用随机梯度下降训练逻辑回归器。分类器与优化器的组合支持利用大量数据在线更新模型；特征投影则使简单分类器也能表示复杂的空间细节。

特征投影与 GP 中核函数的思想相同，但这里使用近似而不是完整协方差。有多种选择，包括 Nystroem \textsuperscript{[1185]}、随机 Fourier 特征（Random Fourier Features）\textsuperscript{[901]} 和稀疏核（Sparse Kernel）\textsuperscript{[755]}，这里采用稀疏核。稀疏核的目标是限制观测产生影响的范围，从而提高收敛速度和计算效率；得到的核函数会在特定距离处严格降为 0。

该核函数通过在待建图空间中按规则间隔放置诱导核，可以将二维或三维空间中的点投影到显著更高的维度。这样就能计算输入数据的高维特征空间表示，并使用小批量随机梯度下降训练逻辑回归分类器。最后，训练通过沿范围测量采样自由空间点，同时将回波加入为障碍物点来完成。

HM 面临的一个挑战是所用核的表达能力。图 5.11 所用的径向基函数（radial basis function，RBF）在二维中是圆、在三维中是球，需要组合多个函数值才能重建环境的复杂细节。因此，核的数量和长度尺度之间存在权衡，二者会影响计算成本、重建细节和插值能力。

\begin{figure}[htbp]
\centering
\includegraphics[width=\linewidth,height=0.78\textheight,keepaspectratio]{5747e487164f38c4984a894030b3612ceeef8b8343c44325b8dcab30ec69c1ad.jpg}
\label{fig:chapter-5-5747e487164f38c4984a894030b3612ceeef8b8343c44325b8dcab30ec69c1ad}
\end{figure}

(a)

\begin{figure}[htbp]
\centering
\includegraphics[width=\linewidth,height=0.78\textheight,keepaspectratio]{c41a8880f986d378ffdd3b2b788bf1040a220b44cc3e9cb23aca49971e288ad2.jpg}
\label{fig:chapter-5-c41a8880f986d378ffdd3b2b788bf1040a220b44cc3e9cb23aca49971e288ad2}
\end{figure}

(b)

\begin{figure}[htbp]
\centering
\includegraphics[width=\linewidth,height=0.78\textheight,keepaspectratio]{94ad4836cb31f41ddbc193757247100d5789dc5d43b9b9ccf2d591faf82cd7ac.jpg}
\label{fig:chapter-5-94ad4836cb31f41ddbc193757247100d5789dc5d43b9b9ccf2d591faf82cd7ac}
\end{figure}

(c)
图 5.11 (a) 机器人使用二维 LiDAR 传感器获得的观测，随后转换为 (b) 自由（绿色）和被占据（红色）点的数据集，再用来训练 (c) 所示的 Hilbert map。

\section{建图中的深度学习}

近年来，利用神经辐射场（neural radiance fields，NeRF，见第 14 章及 \textsuperscript{[764]}）进行新视图合成受到关注；它通过简单的多层感知机（multi-layer perceptron，MLP）实现出色的图像生成结果。一些方法因此研究使用神经表示来估计 SDF。在任意空间位置学习预测 SDF，可以得到连续表示，既能以任意分辨率转换为网格，也能凭借学习函数的插值能力生成更完整的表示。

与隐式表示相似，这种表示从输入数据学习并近似为可在任意位置查询的连续函数。神经表示通常是在给定位姿下离线学习的；近来，人们也开始关注增量方法 \textsuperscript{[1047, 1293]}，以及使用神经表示在线估计位姿的方法 \textsuperscript{[1047, 968]}。

特别地，Sucar 等人 \textsuperscript{[1047]} 的方法基于 RGB-D 帧，利用神经网络预测场景中任意点的 SDF 值。后续方法将特征的空间表示（如体素栅格 \textsuperscript{[862, 1309]}、八叉树 \textsuperscript{[1293]}、点 \textsuperscript{[968, 261]} 等）与神经表示分离。这些方法在空间表示中存储小而有描述力的特征，再利用小型神经网络确定场景中任意点的 SDF 值。这样就将学习函数与空间表示解耦，既可以依靠小型神经解码器将特征转换为有符号距离值，也能够有效应用于室外等大尺度场景。

基于深度学习的建图、重建和 SLAM 领域发展迅速，并结合了经典表示中的思想，例如 surfel splatting \textsuperscript{[548, 740]}，在重建质量和效率方面获得显著成果。第 13 章将介绍 Dust3r 等前馈网络如何从单目视频生成稠密表示。特别是，渲染新视图并在任意位置生成新数据的能力，可能在不依赖仿真环境的情况下用于机器人学习。例如，NeRF 和 Gaussian splatting \textsuperscript{[1318, 103]} 已广受关注，并在多项 SLAM 相关工作中展示出显著潜力；第 14 章将进一步讨论这些方法。

\section{使用考量}

所有地图表示都需要在各自鲜明且往往互补的优点与缺点之间权衡。因此，为给定应用或机器人系统选择地图表示时，必须仔细考虑该地图在所有下游任务中的使用方式。其他需要考虑的因素是运行环境和可用传感器。下面先讨论环境因素，它们会促成一些明确的选择；然后讨论更细致、依赖任务的考量；最后就第 5.4 节所述现有方法给出简要使用建议。

\section{环境因素}

运行环境可以分为结构化和非结构化两类。在自动化工厂等受严格控制的空间中，针对机器人任务及其将遇到的特定对象量身定制的地图表示，在效率和精度方面通常优于通用稠密表示。相反，本章讨论的稠密表示能够对任意形状的物体建模，并可在任何环境中工作。在变化或部分未知的环境中，机器人通常必须能够区分已观测自由空间和未观测空间。该信息使路径规划器能够避免穿越未观测空间的不安全运动，因为这些空间可能被占据；也可用于探索规划。包括点、surfel 和网格在内的显式表面表示，通常无法区分自由空间和未观测空间，而所有基于占据的方法都能做到。基于隐式表面的方法也能提供这种区分，但在许多更专注于重建的应用中，为节省计算和内存成本，远离表面的此类信息经常被丢弃。

另一个考虑因素是可扩展性。显式表示通常比隐式表示更节省内存，因为它们只描述表面本身，并且可以轻易地使保真度适应场景各部分所需的细节。需要自由空间信息时，与固定分辨率体素表示相比，多分辨率方法可以在精度、内存以及捕捉极细物体的能力方面带来显著改进。

最后需要考虑环境中是否包含大量动态物体，以及应在多大程度上对其建模。从地图表示角度看，现有方法大多可归为三类。第一类不考虑动态性，直接将所有测量融合进本章介绍的一种表示。实践中，使用隐式表示时这可能已经足够，因为其自由空间更新通常能很好地消除物体移动后留下的残余。第二类方法试图只将环境静态元素融合到地图中，显式检测并丢弃所有对应动态物体的测量。这类方法在使用显式地图时尤其流行，因为显式地图中的残余更难移除；此类方法通常即使在高度动态的空间中也能生成干净的地图。最后一类方案不仅表示背景，还表示场景中的运动元素。注意，这通常通过混合表示完成，将本章介绍的基本几何表示与对象层级的专用表示结合。关于动态环境中的 SLAM（包括实现上述方法和更先进方法的具体技术）的详细讨论，请参阅第 15 章。

\section{下游任务类型}

除环境外，同样重要的是考虑机器人所需任务需要何种地图信息。任意给定操作通常都能在各种表示上执行，但效率和实现复杂度往往差异很大。最大的差别在于操作沿表面执行，还是在笛卡尔空间中执行。如表 5.2 所示，隐式表示通常可以简单高效地滤波以笛卡尔坐标表达的属性，例如占据状态；相比之下，显式表示适合滤波沿表面表达的属性，例如视觉纹理。这解释了显式表示通常对深度测量质量更敏感，但能够创建非常细致、视觉效果出色的 3D 重建；而隐式方法适合融合 RGB-D 相机数据等噪声深度测量。

就查询而言，显式表示可以直接沿表面遍历，这解释了它们在渲染、图形应用和覆盖路径规划等任务中的流行。不过，要回答笛卡尔坐标中的查询，需要进行最近邻查找等额外步骤。隐式表示恰好相反，因此常用于碰撞检查任务。

\small
\setlength{\tabcolsep}{2pt}
\begin{center}
\begin{tabular}{p{0.317\linewidth}p{0.317\linewidth}p{0.317\linewidth}}
\toprule
操作 & 显式表示中高效 & 隐式表示中高效 \\
\midrule
滤波测量 & 沿表面（纹理等） & 在笛卡尔空间中（占据状态等） \\
查询和遍历 & 在表面坐标中（覆盖规划等） & 在笛卡尔坐标中（碰撞检查等） \\
修改表面 & 几何（形变等） & 拓扑（合并、切割、简化等） \\
\bottomrule
\end{tabular}
\end{center}

表 5.2 显式与隐式表面表示的互补优点和缺点。

最后，显式表示可以高效修改表面的几何形状，包括形变。实践中，地图通常根据不完美且存在漂移的里程计系统提供的位姿估计，通过融合深度测量构建。随着时间推移，累积位姿误差也会导致稠密地图不一致。与 SLAM 系统一样，检测到闭环时可以通过形变稠密地图来消除这些误差。显式和隐式表面都可以形变，但使用显式表示时该操作在本质上更简单、更高效。相反，使用隐式表示可以简化并提高影响表面拓扑或连通性的操作效率。因此，隐式表示常用于合并表面估计、组合或相减物体形状以及简化表面。

还应记住，不同表示可以配合使用，以发挥各自优势。一个很好的混合方法示例是基于 TSDF 的网格化（第 5.4.3 节）：先使用隐式表面表示（TSDF）方便地滤波噪声深度测量，再用 Marching Cubes 将其转换为显式表示（网格）。在判断混合表示的优势是否抵消其带来的开销时，值得考虑如何将转换限制为仅在局部、低频发生。

\section{建图方法总结}

最后总结本章介绍的现有方法的关键差异。先看显式表示：用点集合描述表面简单，所需假设最少，但信息量也最少。除无穷小点外，surfel 在小邻域或片元上表示表面属性。最后，网格显式表示表面连通性，并允许平滑插值其属性。然而，估计表面连通性需要最多假设，有时还伴随显著计算成本。

对隐式表示而言，隐式表面相对于占据地图的一项特殊优势是能提供快速、高质量的距离信息和梯度，有利于基于优化的规划。然而，滤波占据估计所需的假设更少；对于基于体素的方法，占据地图也更擅长捕捉细薄障碍物。在深度测量特别噪声大或稀疏时，基于 GP 和 Hilbert map 的非体素隐式表示能够提供特别好的不确定性估计。由于它们显式考虑几何形状的空间相关性，通常也更擅长插值部分观测到的表面。

学习型方法是一个快速发展的研究领域。对于基于学习的隐式表示，NeRF 已被证明能够支持很有前景的新能力，尤其是语义建模和空间推理。最近，Gaussian splatting \textsuperscript{[552]}（一种与 surfel 相似的显式表示）推动了人们对使用 splatting 的方法 \textsuperscript{[548, 740, 1261]} 的兴趣。研究人员正积极改善这些方法的计算和内存占用，测试其可支持的新能力，并探索如何将其整合到完整机器人系统中。展望未来，我们预计基于学习的方法将通过学习先验，提高稠密表示的通用性和表达能力，并改善其处理噪声测量、不完整观测和动态物体的能力。

\section{致谢}

作者感谢 Lan Wu 在本章准备过程中提供的支持。
